\section{Introduction}
\label{sec:intro}

Albania is one of the most hydropower-dependent electricity systems in Europe. Hydropower made up about
95\% of installed generating capacity at the end of 2018 (2{,}096~MW), and the River Drin cascade, made up of
Fierza (500~MW), Koman (600~MW) and Vau i Dejës (250~MW), is the backbone of domestic supply
\cite{irena2021rra,kesh2026drin}. This structure makes annual electricity supply a function of rainfall and
snowmelt. Reported net imports are closely correlated with annual precipitation \cite{irena2021rra}, and
IRENA's assessment projects that average output from large and small hydropower plants could be roughly 15\% and
20\% lower by 2050 than in 2010 \cite{irena2021rra}. Climate variability is therefore not a remote planning
consideration for Albania. It is the dominant operational risk.

Three developments make the question of \emph{flexibility} timely. First, utility-scale solar photovoltaics
(PV) and batteries have become much cheaper; IRENA reports global weighted-average installed costs of
USD~691/kW for PV, USD~1{,}041/kW for onshore wind and USD~192/kWh for battery storage projects in 2024
\cite{irena2025costs}. Second, the reservoir fleet is itself a large energy store whose seasonal value
depends on how it is coordinated with variable renewables and imports \cite{labadie2004reservoir}. Third,
regional interconnections let a hydro system import in dry periods and export in wet ones. That exposes
the system to regional prices and to the availability of transfer capacity.

This paper asks: \emph{how should a hydropower-dependent electricity system coordinate reservoir operation,
renewable generation, battery storage and electricity imports under uncertain rainfall, fluctuating demand
and financial constraints?} We address eight supporting questions. They cover the reliability effect of
drought, the import-reduction potential of solar, wind and batteries, how storage value changes with drought
severity, budget-constrained investment, the performance of plans optimised for average conditions under
shocks, the value of stochastic and risk-averse planning, the marginal value of each flexibility resource, and
sensitivity to prices, demand, costs and inflows.

\paragraph{Contributions.} The work is an independent research prototype with four contributions:
\begin{enumerate}
\item An open, fully reproducible linear-programming model of the Albanian system. It couples an hourly
energy balance with a physically consistent reservoir water balance (in hm$^3$, with an explicit energy
equivalent), independent battery power and energy sizing, an aggregate interconnector and an unserved-energy
penalty. It is solved with the open-source HiGHS solver and verified against an independent PyPSA formulation.
\item A two-stage stochastic investment programme with a Rockafellar--Uryasev CVaR term. Capacities are
shared first-stage decisions and dispatch is scenario specific at hourly resolution on representative weeks.
Seasonal reservoir coupling is retained through nested annual boundary conditions.
\item A transparent climate-scenario construction that ties drought severities to a reconstructed
2001--2023 inflow index derived from NASA POWER precipitation. Every series is labelled as weather-derived,
reconstructed or synthetic.
\item A systematic marginal-value analysis that compares dual-based and finite-difference values of
battery energy, battery power, PV, wind, import capacity and reservoir volume, and states explicitly when
duals can and cannot be interpreted economically.
\end{enumerate}
The model is embedded in an interactive decision-support application whose numbers come exclusively from
executed optimisation runs. We stress at the outset what the work is not. It is not an operational model,
it has not been validated against observed Albanian grid operations, and its demand and price inputs are
synthetic (Section~\ref{sec:limitations}).
